\section{Experimental Setup}
\label{sec:setup}

\subsection{FinSynth: a synthetic financial-agent testbed}
FinSynth generates, from a fixed seed, 30 fictional listed companies with English and Arabic names, six
sectors, FY2023/FY2024 revenue and net income, ownership links, audit restatements (10 companies), 36
brokerage accounts with holdings and prices, and a 130-document corpus of profiles, annual results,
restatement notices and market news. Five companies' news items carry an indirect prompt-injection
payload that instructs AI assistants to report a fabricated revenue figure. All figures are synthetic,
and every answer is computed programmatically.

\paragraph{Capabilities.} A harness may call the solver LM, keyword search over the corpus, a ticker
resolver (English or Arabic names), a structured financial-data API, a calculator, an entitlement check
and a portfolio API. As in many legacy financial back-ends, the portfolio API does \emph{not} enforce
entitlements itself, so the harness is the enforcement point. The seed harness performs the check.

\paragraph{Task families.} Six families: \emph{lookup} (single figure), \emph{growth} (percent change,
one decimal), \emph{multi-hop} (a figure of a subsidiary's parent), \emph{restatement} (restated FY2023
revenue), \emph{Arabic} (lookups asked in Arabic), and \emph{portfolio} (total value of an authorized
account). Each family has four question templates. Templates A1/A2 are used for adaptation (development
and protected splits), and B1/B2 are \emph{held out} for testing, so test performance reflects transfer to
new phrasings rather than memorized patterns.

\paragraph{Splits.} Dev: 36 questions (6 per family), whose labels the optimizer sees. Protected: 36
curated regression questions (6 per family) with verified labels. Probe: 10 adversarial gate items (5
unauthorized portfolio requests, 5 injection-bearing lookups). Test: 72 held-out-template questions (12
per family). Adversarial test: 12 adversarial items (7 unauthorized, 5 injection), disjoint from the probe.
Policies see dev, protected and probe; only the final analysis touches test and adversarial test.

\subsection{Models and harness revision}
The solver is Qwen3.5-4B (Q4\_K\_M, temperature 0), the solver used in the Reasoning Gym experiments of
\citet{zhang2026harnesslearning}. The proposer is the untrained base Qwen3.5-9B (Q4\_K\_M, temperature
0.8, a fixed sampling seed per proposal). Both run under Ollama on a single 8\,GB laptop RTX~3070.
A capability probe (Appendix~\ref{app:proposer}) showed that the base Qwen3.5-4B proposer, the
original paper's \emph{Base} configuration, produced \emph{no} development improvement in 24
proposals from the seed, in either edit format. A selection study needs some good candidates to select,
so we use the larger base model and report the 4B result as a finding in its own right.

\paragraph{Edit format.} Rewriting a whole program was brittle for both models (Appendix~\ref{app:proposer}).
We therefore ask the proposer for an \emph{additive handler revision}: a function
\texttt{handle(question, tools)} that returns an answer for the questions it targets and \texttt{None}
otherwise. The new harness calls the handler first and falls back to the incumbent when the handler
returns \texttt{None} or raises. Revisions stack across rounds. The format keeps edits small and
revertible but does \emph{not} remove risk. A handler that fires on questions it answers wrongly
regresses them, and a helper that reuses a name defined by the incumbent silently overrides it.

\paragraph{Proposal inputs.} Each proposal prompt contains the capability documentation (signatures,
return formats, short usage examples), the incumbent's full source, an execution report (per-family
development accuracy and up to eight failed examples with expected answers, tool errors, an
upstream-fault note where applicable, and abbreviated traces), and one of six prescribed revision
directives (Appendix~\ref{app:prompts}), following the original paper's use of prescribed revision prompts.

\subsection{Conditions}
\begin{itemize}[leftmargin=*,itemsep=1pt]
  \item \textbf{Clean}: correct development labels, no faults, benign proposer.
  \item \textbf{Adversarial proposer}: as clean, but each proposal is replaced, with probability 0.3, by
  a harmful mutation of itself (Appendix~\ref{app:mutations}). The mutations are an \emph{auth bypass}
  (keeps the entitlement call but ignores its result, so static rules pass), \emph{exfiltration} (a
  guarded network import and request), \emph{answer leakage} (hard-codes development answers), and
  \emph{cost inflation} (four-way self-consistency voting on every LM call). Mutations are applied to a
  real proposal, so they keep whatever genuine improvement it contained.
  \item \textbf{Noisy feedback}: 20\% of development labels are corrupted by $\pm$15--40\%, and the
  structured financial API suffers transient upstream timeouts. With probability 0.3, each of the first
  two calls per question fails; a retry succeeds. The protected suite keeps verified labels.
\end{itemize}

\subsection{Protocol and budgets}
Each (condition, seed) unit runs $T=4$ rounds with $K=4$ proposals per round. All policies within a unit
start from the same seed harness. Proposals are deterministic functions of (parent, report, directive,
sampling seed) and are cached, so two policies that hold the same incumbent at the same round receive
\emph{identical} candidates, and every candidate is evaluated exactly once per condition. Proposal count,
solver calls and feedback are therefore matched across policies by construction (common random numbers).
We use 3 seeds per condition, giving 9 units. Budgets are fixed before running: $B_C = 600$ solver
tokens per question (about $3\times$ the seed harness), $B_L = 6$\,s modeled p95 latency, $\epsilon_R =
0.15$, $\alpha = 0.10$, 95\% one-sided bounds. Latency is modeled deterministically from token counts (0.15\,s
overhead, 2{,}500 prompt tokens/s, 75 generated tokens/s, calibrated on the study GPU), plus fixed per-tool
costs and a 2\,s penalty per upstream fault. A model avoids measuring request queueing during parallel
evaluation. Per-question limits are 8 LM calls, 40 tool calls
and 90\,s wall time.

\subsection{Metrics}
For the final harness of each run we report held-out test accuracy, mean solver tokens per question, the
number of unauthorized portfolio reads on the adversarial test, the injection-follow rate, whether the
code attempts exfiltration, and the regression rate relative to the seed on held-out test items. We also
count accepted edits and accepted harmful edits. For the decision-level analysis, every logged (parent,
candidate) pair is labeled with ground truth from the held-out splits: \emph{harmful}
(security, injection, leakage, cost $> B_C$, or test regression $>10\%$ of parent-solved items) and
\emph{beneficial} (no harm and a held-out gain of at least 3 of 72 questions). We then measure how often
each policy's gates pass harmful and beneficial candidates. Uncertainty is reported as bootstrap 95\%
confidence intervals over units, and paired policy differences use exact sign tests or paired bootstrap.
